% Generated draft: requires mathematical and source review.
\documentclass[11pt]{article}
\usepackage[margin=25mm]{geometry}
\usepackage{fontspec}
\setmainfont{Latin Modern Roman}
\usepackage{amsmath,amssymb,mathtools,unicode-math}
\setmathfont{Latin Modern Math}
\usepackage{enumitem,needspace,xurl,hyperref}
\setlength{\parindent}{0pt}
\setlength{\parskip}{5pt}
\setlength{\emergencystretch}{3em}
\newcommand{\N}{\mathbb{N}}
\newcommand{\Z}{\mathbb{Z}}
\newcommand{\Q}{\mathbb{Q}}
\newcommand{\R}{\mathbb{R}}
\newcommand{\C}{\mathbb{C}}
\newcommand{\F}{\mathbb{F}}
\newcommand{\A}{\mathbb{A}}
\newcommand{\PP}{\mathbb{P}}
\newcommand{\E}{\mathbb{E}}
\newcommand{\eps}{\varepsilon}
\newcommand{\dd}{\,\mathrm{d}}
\begin{document}
\section*{UnsolvedMath ID: 30006317}
\textbf{Source problem code:} catalogue entry
\subsection{Definitions and mathematical statement}
The supplied record describes a conjectural framework whose objects and hypotheses are not fully defined. Let \(G\) denote the group appearing in that framework and \(G^\vee\) its intended Langlands dual. The base field, admissible class of groups, geometric category, and precise meaning of a Hamiltonian \(G\)-space are not specified. Write \(\mathcal A_G\) and \(\mathcal B_G\) for the two theories named in the source; these symbols do not constitute definitions of those theories.

\textbf{Conjectural existence and compatibility problem.} For a suitable class of Hamiltonian \(G\)-spaces \(M\), construct boundary conditions associated with \(M\) in both \(\mathcal A_G\) and \(\mathcal B_G\), and construct a Hamiltonian \(G^\vee\)-space \(M^\vee\) such that a local geometric Langlands equivalence
\[
\Phi_G:\mathcal A_G\simeq\mathcal B_{G^\vee}
\]
matches the boundary condition associated with \(M\) on the \(\mathcal A_G\) side with that associated with \(M^\vee\) on the \(\mathcal B_{G^\vee}\) side. The original excerpt says this is expected for ``many'' Hamiltonian spaces; it neither specifies that class nor asserts the construction for every Hamiltonian \(G\)-space. The formulation is therefore a conditional research problem, not a fully quantified conjecture with supplied hypotheses.

\textbf{Cotangent and unramified formulation.} Let \(M=T^*X\) be the cotangent case in the source, with a \(G\)-action on \(X\) implicit in the expressions below. Let \(\mathcal D\) denote the formal disc and \(\mathring{\mathcal D}\) its punctured disc over the unspecified base. Here \(G(\mathcal D)\), \(G(\mathring{\mathcal D})\), and \(X(\mathring{\mathcal D})\) denote the corresponding spaces of points or maps, in a geometric setting still to be specified. Write \(\mathcal B^A_M\) as a placeholder for the boundary condition associated with \(M\) in \(\mathcal A_G\). The source proposes the informal identification
\[
\mathcal B^A_M(\mathring{\mathcal D})
\simeq \operatorname{Shv}\bigl(X(\mathring{\mathcal D})\bigr)
\quad\text{in}\quad
\mathcal A_G(\mathring{\mathcal D})
=\operatorname{Shv}\bigl(G(\mathring{\mathcal D})\bigr)\text{-}\operatorname{modcat}.
\]
The notation on the right denotes a category of module categories over the intended sheaf category with its intended monoidal structure. Neither this sheaf theory nor the required categorical foundations are defined in general in the supplied excerpt.

The more accessible, but still incompletely founded, unramified expression in the source is
\[
\begin{aligned}
\bigl\langle\mathcal A_G(\mathcal D),\mathcal B^A_M(\mathcal D)\bigr\rangle
&\simeq
\operatorname{Shv}\bigl(G(\mathcal D)\backslash G(\mathring{\mathcal D})\bigr)
\underset{\operatorname{Shv}(G(\mathring{\mathcal D}))}{\otimes}
\operatorname{Shv}\bigl(X(\mathring{\mathcal D})\bigr)\\
&\simeq
\operatorname{Shv}_{G(\mathcal D)}\bigl(X(\mathring{\mathcal D})\bigr).
\end{aligned}
\]
The last term is intended to mean \(G(\mathcal D)\)-equivariant sheaves. The pairing, quotient, relative tensor product, sheaf theory, and module structures require compatible definitions. The task is to make these constructions precise and relate them to the proposed matching of dual boundary conditions. The excerpt supplies no explicit category on the \(M^\vee\) side for this unramified expression, so no specific equivalence with such a category can be stated from the record alone.
\subsection{Short English statement}
For an unspecified suitable class of Hamiltonian spaces, construct dual spaces whose boundary conditions correspond under local geometric Langlands duality. In the cotangent case, the source informally expresses the unramified construction using equivariant sheaves on punctured-disc maps into \(X\). Its general foundations and exact hypotheses remain unspecified, and the original excerpt does not claim the construction for every Hamiltonian space.
\subsection{Sources}
\begin{itemize}
\item Supplied local UnsolvedMath registry record: permanent ID \(30006317\), source problem code \texttt{catalogue entry}, title \emph{The Local Geometric Langlands Conjecture}. The record gives proposed year 2025 and names no individual proposer.
\item \emph{Arbeitsgemeinschaft: Relative Langlands Duality} (2025), Oberwolfach Reports, pp.\ 811--882, DOI \texttt{10.4171/owr/2025/18}; supplied excerpt headed ``2.2. The Local Geometric Conjecture.'' \url{https://doi.org/10.4171/owr/2025/18}.
\item Ben-Zvi, Sakellaridis, and Venkatesh, \emph{Relative Langlands Duality}, supplied literature reference: \url{https://www.math.ias.edu/~akshay/research/BZSVpaperV1.pdf}.
\item Devalapurkar, \emph{Derived geometric Satake for PGL products}, supplied literature-reference label: \url{https://sanathdevalapurkar.github.io/files/PGL2-cubes.pdf}.
\item Supplied literature assessment dated \texttt{2026-08-22}: the record labels the problem open, reports that the extracted report lacks the general technical foundations, and describes later results as proving special rank-one or group cases rather than the stated general equivalence. This is supplied dated curation, not an independently verified assessment; its time-sensitive status and cited claims require independent checking.
\end{itemize}
\end{document}
